\documentclass[12pt,a4paper]{article}

% ─── Encodage & langue ───────────────────────────────────────────────────────
\usepackage[utf8]{inputenc}
\usepackage[T1]{fontenc}
\usepackage[french]{babel}

% ─── Mise en page ────────────────────────────────────────────────────────────
\usepackage[top=2.5cm, bottom=2.5cm, left=2.5cm, right=2.5cm]{geometry}
\usepackage{setspace}
\onehalfspacing

% ─── Mathématiques ───────────────────────────────────────────────────────────
\usepackage{amsmath, amssymb, amsthm}

% ─── Tableaux & figures ──────────────────────────────────────────────────────
\usepackage{booktabs}
\usepackage{array}
\usepackage{multirow}
\usepackage{float}
\usepackage{graphicx}
\usepackage{caption}

% ─── Code source ─────────────────────────────────────────────────────────────
\usepackage{listings}
\usepackage{xcolor}

\definecolor{codebg}{RGB}{248,248,248}
\definecolor{codecomment}{RGB}{100,100,100}
\definecolor{codekeyword}{RGB}{0,0,180}
\definecolor{codestring}{RGB}{163,21,21}
\definecolor{codenumber}{RGB}{128,128,128}
\definecolor{codeframe}{RGB}{200,200,200}

\lstdefinestyle{cppstyle}{
  language=C++,
  backgroundcolor=\color{codebg},
  basicstyle=\ttfamily\small,
  keywordstyle=\color{codekeyword}\bfseries,
  commentstyle=\color{codecomment}\itshape,
  stringstyle=\color{codestring},
  numberstyle=\tiny\color{codenumber},
  numbers=left,
  numbersep=8pt,
  frame=single,
  rulecolor=\color{codeframe},
  breaklines=true,
  breakatwhitespace=true,
  tabsize=4,
  showstringspaces=false,
  captionpos=b,
  xleftmargin=15pt,
  morekeywords={vector, string, numeric_limits}
}

\lstdefinestyle{bashstyle}{
  language=bash,
  backgroundcolor=\color{codebg},
  basicstyle=\ttfamily\small,
  keywordstyle=\color{codekeyword}\bfseries,
  commentstyle=\color{codecomment}\itshape,
  stringstyle=\color{codestring},
  frame=single,
  rulecolor=\color{codeframe},
  breaklines=true,
  showstringspaces=false,
  captionpos=b,
  xleftmargin=15pt,
}

\lstdefinestyle{outputstyle}{
  backgroundcolor=\color{codebg},
  basicstyle=\ttfamily\small,
  frame=single,
  rulecolor=\color{codeframe},
  breaklines=true,
  captionpos=b,
  xleftmargin=15pt,
  numbers=none,
}

% ─── Couleurs & boîtes ───────────────────────────────────────────────────────
\usepackage{tcolorbox}
\tcbuselibrary{skins, breakable}

\newtcolorbox{resultbox}{
  colback=green!5!white,
  colframe=green!60!black,
  fonttitle=\bfseries,
  title=Résultat,
  breakable
}

\newtcolorbox{notebox}{
  colback=blue!5!white,
  colframe=blue!50!black,
  fonttitle=\bfseries,
  title=Remarque,
  breakable
}

\newtcolorbox{formulebox}{
  colback=orange!5!white,
  colframe=orange!70!black,
  fonttitle=\bfseries,
  breakable
}

% ─── Liens ───────────────────────────────────────────────────────────────────
\usepackage[colorlinks=true, linkcolor=blue!70!black, urlcolor=blue]{hyperref}

% ─── En-tête / pied de page ──────────────────────────────────────────────────
\usepackage{fancyhdr}
\pagestyle{fancy}
\fancyhf{}
\lhead{\textit{Master informatique — parcours IA}}
\rhead{\textit{Modèles stochastiques}}
\cfoot{\thepage}
\renewcommand{\headrulewidth}{0.4pt}

% ─── Titre ───────────────────────────────────────────────────────────────────
\title{
  \vspace{-1cm}
  \rule{\linewidth}{1.5pt}\\[0.4cm]
  {\LARGE \textbf{TP 2 : Modèles de Markov cachés}}\\[0.3cm]
  {\large Rapport de travaux pratiques}\\[0.2cm]
  \rule{\linewidth}{1pt}
}
\author{Master Informatique --- Parcours Intelligence Artificielle\\
        Avignon Université}
\date{Mars 2025}

% ─────────────────────────────────────────────────────────────────────────────
\begin{document}

\maketitle
\thispagestyle{empty}
\tableofcontents
\newpage

% ═════════════════════════════════════════════════════════════════════════════
\section{Prédiction de la meilleure séquence d'états}
% ═════════════════════════════════════════════════════════════════════════════

\subsection{Implémentation de l'algorithme de Viterbi}

L'algorithme de Viterbi recherche la séquence d'états cachés $Q^* = (q_1^*, \ldots, q_T^*)$
qui maximise la probabilité jointe $P(Q, O \mid \lambda)$. Tous les calculs sont effectués
en log-probabilités pour éviter les sous-débordements numériques.

\subsubsection*{Formulation}

On définit la variable $\delta_t(i)$ comme le log de la probabilité maximale d'une
séquence d'états se terminant en $s_i$ au temps $t$ et ayant généré $o_1, \ldots, o_t$ :

\begin{formulebox}
\textbf{Initialisation} ($t = 0$) :
\[
  \delta_0(i) = \log \pi_i + \log B_{i,o_0}, \quad \psi_0(i) = 0
\]

\textbf{Récurrence} ($t = 1, \ldots, T-1$) :
\[
  \delta_t(j) = \max_{i}\bigl[\delta_{t-1}(i) + \log A_{ij}\bigr] + \log B_{j,o_t}
\]
\[
  \psi_t(j) = \arg\max_{i}\bigl[\delta_{t-1}(i) + \log A_{ij}\bigr]
\]

\textbf{Terminaison} :
\[
  q^*_{T-1} = \arg\max_{i}\, \delta_{T-1}(i)
\]

\textbf{Rétropiste} ($t = T-2, \ldots, 0$) :
\[
  q^*_t = \psi_{t+1}(q^*_{t+1})
\]
\end{formulebox}

\subsubsection*{Code C++ implémenté}

\begin{lstlisting}[style=cppstyle, caption={Algorithme de Viterbi --- \texttt{model.cpp}}]
vector<int>* Model::viterbi(vector<int>* seq_obs)
{
    int T = seq_obs->size();
    double** delta = new double*[T];
    int**    psi   = new int*[T];
    for (int t = 0; t < T; t++) {
        delta[t] = new double[n];
        psi[t]   = new int[n];
    }
    // Initialisation
    for (int i = 0; i < n; i++) {
        delta[0][i] = pi[i] + B[i][(*seq_obs)[0]];
        psi[0][i]   = 0;
    }
    // Recurrence
    for (int t = 1; t < T; t++) {
        for (int j = 0; j < n; j++) {
            double best = -numeric_limits<double>::infinity();
            int best_i = 0;
            for (int i = 0; i < n; i++) {
                double val = delta[t-1][i] + A[i][j];
                if (val > best) { best = val; best_i = i; }
            }
            delta[t][j] = best + B[j][(*seq_obs)[t]];
            psi[t][j]   = best_i;
        }
    }
    // Terminaison
    double best = -numeric_limits<double>::infinity();
    int best_state = 0;
    for (int i = 0; i < n; i++)
        if (delta[T-1][i] > best) { best = delta[T-1][i]; best_state = i; }
    // Retropiste
    vector<int>* res = new vector<int>(T);
    (*res)[T-1] = best_state;
    for (int t = T-2; t >= 0; t--)
        (*res)[t] = psi[t+1][(*res)[t+1]];
    // Nettoyage memoire
    for (int t = 0; t < T; t++) { delete[] delta[t]; delete[] psi[t]; }
    delete[] delta; delete[] psi;
    return res;
}
\end{lstlisting}

\subsection{Test sur les données météo}

\subsubsection*{Commandes exécutées}

\begin{lstlisting}[style=bashstyle]
./use_hmm test ../data/weather/weather.true.model \
               ../data/weather/weather-test.csv \
               /tmp/weather-pred.txt
python3 compErrors.py ../data/weather/weather-test.ref /tmp/weather-pred.txt
\end{lstlisting}

\subsubsection*{Résultat}

\begin{lstlisting}[style=outputstyle]
Errors: 0/18 (0.00%)
-------------------------
|     ERRORS BY STATE    |
| chaud:    0 (000.00%) |
| doux :    0 (000.00%) |
| froid:    0 (000.00%) |
-------------------------
\end{lstlisting}

\begin{resultbox}
Les prédictions de Viterbi sont \textbf{identiques} au fichier de référence
\texttt{weather-test.ref} : \textbf{0\% d'erreurs} sur les 18 observations de test.
L'implémentation est correcte.
\end{resultbox}

% ═════════════════════════════════════════════════════════════════════════════
\section{Application en bio-informatique}
% ═════════════════════════════════════════════════════════════════════════════

\subsection{Modèle HMM pour la détection de régions CG}

Le modèle HMM à deux états est défini directement d'après les paramètres du sujet.

\begin{table}[H]
\centering
\caption{Matrice de transition $A$}
\begin{tabular}{lcc}
\toprule
 & \textbf{haut taux GC (+)} & \textbf{bas taux GC (−)} \\
\midrule
\textbf{haut taux GC (+)} & 0,99 & 0,01 \\
\textbf{bas taux GC (−)}  & 0,01 & 0,99 \\
\bottomrule
\end{tabular}
\end{table}

\begin{table}[H]
\centering
\caption{Matrice d'émission $B$ et probabilités initiales $\pi$}
\begin{tabular}{lcccc}
\toprule
 & $v_t = A$ & $v_t = G$ & $v_t = C$ & $v_t = T$ \\
\midrule
$P(o_t = v_t \mid q_t = +)$ & 0,20 & 0,30 & 0,30 & 0,20 \\
$P(o_t = v_t \mid q_t = -)$ & 0,30 & 0,20 & 0,20 & 0,30 \\
\bottomrule
\end{tabular}
\end{table}

Probabilités initiales : $\pi = (0{,}5 \;;\; 0{,}5)$.

\subsubsection*{Commandes exécutées}

\begin{lstlisting}[style=bashstyle]
./use_hmm test ../data/genetics/CG-region.model \
               ../data/genetics/CG-region.csv \
               /tmp/CG-pred.txt
python3 compErrors.py ../data/genetics/CG-region.ref /tmp/CG-pred.txt
\end{lstlisting}

\subsubsection*{Résultat}

\begin{lstlisting}[style=outputstyle]
Errors: 16685/100000 (16.68%)
-------------------------
|     ERRORS BY STATE    |
| +    : 8250 (016.64%) |
| -    : 8435 (016.73%) |
-------------------------
\end{lstlisting}

\begin{resultbox}
\textbf{Précision du HMM : 83,32\%} (taux d'erreur : 16,68\% sur 100\,000 bases).
\end{resultbox}

\begin{notebox}
Le taux d'erreur est quasi-symétrique entre les deux états ($\approx 16{,}7\%$ chacun).
Cela s'explique par le fait que les distributions d'émission des deux états sont proches :
les régions à haut taux de GC émettent G/C à 60\% et A/T à 40\%, tandis que les régions
à bas taux font l'inverse. Cette faible discrimination rend difficile la séparation exacte
des deux types de régions, d'autant que les transitions sont identiques (0,99/0,01)
pour les deux états.
\end{notebox}

% ═════════════════════════════════════════════════════════════════════════════
\section{Apprentissage par maximum de vraisemblance}
% ═════════════════════════════════════════════════════════════════════════════

\subsection{Formules MLE}

Étant donné un corpus supervisé de $N$ séquences $\{(O^{(k)}, Q^{(k)})\}_{k=1}^{N}$,
où chaque observation $o_t^{(k)}$ est annotée par son état caché $q_t^{(k)}$, les
estimateurs du maximum de vraisemblance sont :

\begin{formulebox}
\textbf{Probabilités initiales} :
\[
  \hat{\pi}_i = \frac{\#\{\text{séquences commençant par l'état } s_i\}}{N}
             = \frac{\displaystyle\sum_{k=1}^{N} \mathbf{1}\bigl[q_1^{(k)} = s_i\bigr]}{N}
\]

\textbf{Probabilités de transition} :
\[
  \hat{a}_{ij} =
  \frac{\displaystyle\sum_{k=1}^{N}\sum_{t=1}^{T_k - 1}
        \mathbf{1}\bigl[q_t^{(k)}=s_i,\; q_{t+1}^{(k)}=s_j\bigr]}
       {\displaystyle\sum_{k=1}^{N}\sum_{t=1}^{T_k - 1}
        \mathbf{1}\bigl[q_t^{(k)}=s_i\bigr]}
\]

\textbf{Probabilités d'émission} :
\[
  \hat{b}_{jk} =
  \frac{\displaystyle\sum_{k=1}^{N}\sum_{t=1}^{T_k}
        \mathbf{1}\bigl[q_t^{(k)}=s_j,\; o_t^{(k)}=v_k\bigr]}
       {\displaystyle\sum_{k=1}^{N}\sum_{t=1}^{T_k}
        \mathbf{1}\bigl[q_t^{(k)}=s_j\bigr]}
\]
\end{formulebox}

Ces estimateurs maximisent la log-vraisemblance supervisée
$\log P(O, Q \mid \lambda) = \sum_k \log P(O^{(k)}, Q^{(k)} \mid \lambda)$.

\subsection{Implémentation de \texttt{train\_MLE}}

\begin{lstlisting}[style=cppstyle, caption={Apprentissage MLE supervisé --- \texttt{model.cpp}}]
void Model::train_MLE(vector<vector<int>*>* data_obs,
                      vector<vector<int>*>* data_state, string filename)
{
    // Initialisation des compteurs
    int* count_start  = new int[n]();
    int* count_unigram = new int[n]();
    int** count_bigram = new int*[n];
    int** count_state_obs = new int*[n];
    for (int i = 0; i < n; i++) {
        count_bigram[i]    = new int[n]();
        count_state_obs[i] = new int[M]();
    }
    // Comptage
    int N = data_obs->size();
    for (int k = 0; k < N; k++) {
        vector<int>* obs   = (*data_obs)[k];
        vector<int>* state = (*data_state)[k];
        int T = obs->size();
        count_start[(*state)[0]]++;
        for (int t = 0; t < T; t++) {
            int s = (*state)[t], o = (*obs)[t];
            count_state_obs[s][o]++;
            if (t < T-1) {
                count_bigram[s][(*state)[t+1]]++;
                count_unigram[s]++;
            }
        }
    }
    // Conversion en log-probabilites
    for (int i = 0; i < n; i++)
        pi[i] = (count_start[i] > 0)
                ? log((double)count_start[i] / N) : LOG_MINUS_INFTY;
    for (int i = 0; i < n; i++)
        for (int j = 0; j < n; j++)
            A[i][j] = (count_bigram[i][j] > 0)
                      ? log((double)count_bigram[i][j]/count_unigram[i])
                      : LOG_MINUS_INFTY;
    for (int j = 0; j < n; j++) {
        int total = 0;
        for (int k = 0; k < M; k++) total += count_state_obs[j][k];
        for (int k = 0; k < M; k++)
            B[j][k] = (count_state_obs[j][k] > 0)
                      ? log((double)count_state_obs[j][k] / total)
                      : LOG_MINUS_INFTY;
    }
    init_param_EOS();
    save_hmm(filename);
    // Nettoyage memoire ...
}
\end{lstlisting}

\subsection{Test sur les données météo}

\subsubsection*{Commandes exécutées}

\begin{lstlisting}[style=bashstyle]
./use_hmm train_sup /tmp/weather-mle.model ../data/weather/weather-500seq.csv
./use_hmm test /tmp/weather-mle.model ../data/weather/weather-test.csv \
               /tmp/weather-mle-pred.txt
python3 compErrors.py ../data/weather/weather-test.ref /tmp/weather-mle-pred.txt
\end{lstlisting}

\subsubsection*{Paramètres appris (extrait)}

\begin{table}[H]
\centering
\caption{Matrice d'émission $\hat{B}$ apprise (à comparer avec le vrai modèle)}
\begin{tabular}{lccc}
\toprule
\textbf{État} & \textbf{été} & \textbf{mi-saison} & \textbf{hiver} \\
\midrule
chaud & 0,799 & 0,201 & $\approx 0$ \\
doux  & 0,399 & 0,394 & 0,208 \\
froid & 0,110 & 0,111 & 0,779 \\
\midrule
\textit{vrai modèle} & & & \\
chaud & 0,800 & 0,190 & 0,010 \\
doux  & 0,500 & 0,400 & 0,100 \\
froid & 0,010 & 0,200 & 0,790 \\
\bottomrule
\end{tabular}
\end{table}

\subsubsection*{Résultat}

\begin{lstlisting}[style=outputstyle]
Errors: 0/18 (0.00%)
-------------------------
|     ERRORS BY STATE    |
| chaud:    0 (000.00%) |
| doux :    0 (000.00%) |
| froid:    0 (000.00%) |
-------------------------
\end{lstlisting}

\begin{resultbox}
Le modèle appris par MLE sur 500 séquences atteint \textbf{0\% d'erreurs} sur les données
de test, identique au vrai modèle. Les paramètres estimés sont très proches des paramètres
du vrai modèle, ce qui valide l'implémentation de \texttt{train\_MLE}.
\end{resultbox}

% ═════════════════════════════════════════════════════════════════════════════
\section{Application en traitement automatique des langues}
% ═════════════════════════════════════════════════════════════════════════════

\subsection{Traitement des mots inconnus}

À partir de l'examen des fonctions \texttt{read\_vocab} et \texttt{read\_sequences},
voici comment le programme gère les mots inconnus :

\paragraph{Lors de l'apprentissage --- \texttt{read\_vocab} :}
Lors d'un premier parcours du fichier d'entraînement, chaque mot est compté.
Tout mot dont le compte est \textbf{strictement inférieur au seuil}
\texttt{UNK\_THRESHOLD = 2} est retiré du vocabulaire et ses occurrences sont
attribuées au token spécial \texttt{"UNK"}.
Ce token devient une observation à part entière du modèle, avec ses propres
probabilités d'émission $B[\text{état}][\texttt{UNK}]$ apprises sur l'ensemble
des mots rares du corpus d'entraînement.

\paragraph{Lors du test --- \texttt{read\_sequences} (train=false) :}
Si un mot du fichier de test est \textbf{absent du dictionnaire} \texttt{dic\_obs}
du modèle (jamais vu ou trop rare à l'entraînement), il est remplacé par l'index
de \texttt{UNK} dans la séquence d'observations.
Sa forme originale est conservée dans une \textbf{file FIFO} \texttt{unk\_obs}
pour être restituée à l'écriture.

\paragraph{Lors de l'écriture --- \texttt{write\_predict} :}
Si l'observation courante correspond à l'index de \texttt{UNK}, le mot original
est extrait depuis la file \texttt{unk\_obs} et affiché à la place.
L'étiquette PoS attribuée est celle calculée par Viterbi pour le token \texttt{UNK}.

\begin{notebox}
Contrairement aux applications météo et bio-informatique où le vocabulaire
des observations est \textbf{fermé} (le nombre de symboles possibles est fixe),
le vocabulaire des mots en traitement des langues est \textbf{ouvert} : de nouveaux
mots peuvent apparaître au test. Le mécanisme \texttt{UNK} permet de gérer ce cas
en regroupant tous les mots inconnus sous une catégorie générique.
\end{notebox}

\subsection{Résultats du POS-tagging}

\subsubsection*{Entraînement}

\begin{lstlisting}[style=bashstyle]
./use_hmm train_sup /tmp/pos-mle.model \
          ../data/POS-tagging/fr-train.word-pos EOS
\end{lstlisting}

Modèle appris : \textbf{18 états} (catégories PoS + EOS), \textbf{19\,070 observations}.

\subsubsection*{Tableau synthétique des erreurs}

\begin{table}[H]
\centering
\caption{Taux d'erreurs par corpus de test}
\begin{tabular}{lrrr}
\toprule
\textbf{Corpus de test} & \textbf{Erreurs} & \textbf{Total} & \textbf{Taux d'erreur} \\
\midrule
\texttt{fr-test} (tous sous-corpus) & 6\,604 & 81\,569 & \textbf{8,10\%} \\
\texttt{UD\_French-FQB}             & 2\,046 & 24\,135 & 8,48\% \\
\texttt{UD\_French-GSD}             &    679 & 10\,019 & 6,78\% \\
\texttt{UD\_French-ParTUT}          &    168 &  2\,603 & 6,45\% \\
\texttt{UD\_French-PUD}             & 2\,200 & 24\,735 & 8,89\% \\
\texttt{UD\_French-Sequoia}         &    628 & 10\,048 & 6,25\% \\
\texttt{UD\_French-Spoken}          &    883 & 10\,029 & 8,80\% \\
\bottomrule
\end{tabular}
\end{table}

\subsubsection*{Étiquettes les plus difficiles à prédire}

\begin{table}[H]
\centering
\caption{Taux d'erreur par étiquette PoS sur \texttt{fr-test}}
\begin{tabular}{lccc}
\toprule
\textbf{Étiquette} & \textbf{Description} & \textbf{Erreurs} & \textbf{Taux} \\
\midrule
\texttt{X}     & Mots étrangers / inclassables & 164   & \textbf{46,86\%} \\
\texttt{ADJ}   & Adjectifs                     & 1\,213 & \textbf{24,05\%} \\
\texttt{ADV}   & Adverbes                      &   577  & 15,35\% \\
\texttt{VERB}  & Verbes                        & 1\,052 & 14,71\% \\
\texttt{NUM}   & Numéraux                      &   164  & 13,53\% \\
\texttt{SYM}   & Symboles                      &    10  & 13,16\% \\
\texttt{PROPN} & Noms propres                  &   597  & 12,97\% \\
\texttt{NOUN}  & Noms communs                  & 1\,609 & 10,73\% \\
\texttt{SCONJ} & Conj. de subordination        &    55  & 5,84\% \\
\texttt{PRON}  & Pronoms                       &   343  & 6,35\% \\
\texttt{AUX}   & Auxiliaires                   &    53  & 1,31\% \\
\texttt{PUNCT} & Ponctuation                   &     1  & 0,01\% \\
\bottomrule
\end{tabular}
\end{table}

\paragraph{Analyse des difficultés :}
\begin{itemize}
  \item \textbf{X (46,9\%)} : mots hors-norme (emprunts, abréviations, erreurs)
        souvent inconnus du modèle → traités comme \texttt{UNK}, catégorie très minoritaire.
  \item \textbf{ADJ (24,1\%)} : forte ambiguïté morphologique en français.
        Un participe passé peut être ADJ ou VERB selon le contexte
        (ex. \textit{ouvert}, \textit{construit}).
  \item \textbf{ADV (15,4\%)} et \textbf{VERB (14,7\%)} : nombreux mots
        polysémiques dont l'étiquette dépend du contexte à longue portée,
        que le HMM (Markov d'ordre 1) ne capture pas.
  \item \textbf{PROPN (13,0\%)} : les noms propres sont fréquemment inconnus
        du vocabulaire d'entraînement.
  \item \textbf{PUNCT (0,01\%)} et \textbf{AUX (1,31\%)} : catégories
        très régulières, quasi-injectives mots $\to$ étiquette.
\end{itemize}

% ═════════════════════════════════════════════════════════════════════════════
\section{Calcul de scores de confiance}
% ═════════════════════════════════════════════════════════════════════════════

\subsection{Fonctions \texttt{alpha} et \texttt{beta}}

\subsubsection*{Forward --- $\alpha_t(i) = \log P(o_1, \ldots, o_t,\; q_t = s_i)$}

\begin{formulebox}
\textbf{Initialisation} :
\[
  \alpha_0(i) = \log \pi_i + \log B_{i,o_0}
\]
\textbf{Récurrence} :
\[
  \alpha_t(j) = \ln\!\left(\sum_{i=1}^{n} e^{\alpha_{t-1}(i) + \log A_{ij}}\right) + \log B_{j,o_t}
\]
La somme est calculée via \texttt{ln\_of\_sum} pour éviter les sous-débordements.
\end{formulebox}

\begin{lstlisting}[style=cppstyle, caption={\texttt{alpha} --- \texttt{model.cpp}}]
double** Model::alpha(vector<int>* seq_obs) {
    int T = seq_obs->size();
    double** res = new double*[T];
    for (int t = 0; t < T; t++) res[t] = new double[n];
    // Initialisation
    for (int i = 0; i < n; i++)
        res[0][i] = pi[i] + B[i][(*seq_obs)[0]];
    // Recurrence
    for (int t = 1; t < T; t++) {
        for (int j = 0; j < n; j++) {
            double* tmp = new double[n];
            for (int i = 0; i < n; i++)
                tmp[i] = res[t-1][i] + A[i][j];
            res[t][j] = ln_of_sum(tmp, 0, n-1) + B[j][(*seq_obs)[t]];
            delete[] tmp;
        }
    }
    return res;
}
\end{lstlisting}

\subsubsection*{Backward --- $\beta_t(i) = \log P(o_{t+1}, \ldots, o_T \mid q_t = s_i)$}

\begin{formulebox}
\textbf{Initialisation} :
\[
  \beta_{T-1}(i) = \log(1) = 0
\]
\textbf{Récurrence} (de $T-2$ à $0$) :
\[
  \beta_t(i) = \ln\!\left(\sum_{j=1}^{n} e^{\log A_{ij} + \log B_{j,o_{t+1}} + \beta_{t+1}(j)}\right)
\]
\end{formulebox}

\begin{lstlisting}[style=cppstyle, caption={\texttt{beta} --- \texttt{model.cpp}}]
double** Model::beta(vector<int>* seq_obs) {
    int T = seq_obs->size();
    double** res = new double*[T];
    for (int t = 0; t < T; t++) res[t] = new double[n];
    // Initialisation
    for (int i = 0; i < n; i++) res[T-1][i] = 0.0;
    // Recurrence
    for (int t = T-2; t >= 0; t--) {
        for (int i = 0; i < n; i++) {
            double* tmp = new double[n];
            for (int j = 0; j < n; j++)
                tmp[j] = A[i][j] + B[j][(*seq_obs)[t+1]] + res[t+1][j];
            res[t][i] = ln_of_sum(tmp, 0, n-1);
            delete[] tmp;
        }
    }
    return res;
}
\end{lstlisting}

\subsection{Fonction \texttt{pgen}}

\begin{formulebox}
\[
  \log P(O^k \mid \mathcal{M}) = \ln\!\left(\sum_{i=1}^{n} e^{\alpha_T(i)}\right)
  = \texttt{ln\_of\_sum}(\alpha[T-1],\; 0,\; n-1)
\]
\end{formulebox}

\begin{lstlisting}[style=cppstyle, caption={\texttt{pgen} --- \texttt{model.cpp}}]
double Model::pgen(vector<int>* seq_obs, double** alpha) {
    int T = seq_obs->size();
    return ln_of_sum(alpha[T-1], 0, n-1);
}
\end{lstlisting}

\subsection{Fonction \texttt{posterior}}

La probabilité a posteriori que l'état $s_i$ ait émis l'observation au temps $t$ est :

\begin{formulebox}
\[
  \gamma_t(i) = \frac{\alpha_t(i)\,\beta_t(i)}{P(O^k \mid \mathcal{M})}
  \implies
  \log \gamma_t(s_t) = \alpha_t(s_t) + \beta_t(s_t) - \log P(O^k \mid \mathcal{M})
\]
où $s_t$ est l'état prédit par Viterbi au temps $t$.
\end{formulebox}

\begin{lstlisting}[style=cppstyle, caption={\texttt{posterior} --- \texttt{model.cpp}}]
vector<double>* Model::posterior(vector<int>* seq_obs, vector<int>* seq_state) {
    int T = seq_obs->size();
    double** alpha_mat = Model::alpha(seq_obs);
    double** beta_mat  = Model::beta(seq_obs);
    double denom = Model::pgen(seq_obs, alpha_mat);
    vector<double>* res = new vector<double>(T);
    for (int t = 0; t < T; t++) {
        int s = (*seq_state)[t];
        (*res)[t] = alpha_mat[t][s] + beta_mat[t][s] - denom;
    }
    for (int t = 0; t < T; t++) {
        delete[] alpha_mat[t]; delete[] beta_mat[t];
    }
    delete[] alpha_mat; delete[] beta_mat;
    return res;
}
\end{lstlisting}

\subsection{Vérification et affichage}

\subsubsection*{Vérification sur les données météo}

\begin{lstlisting}[style=bashstyle]
./use_hmm test ../data/weather/weather.true.model \
               ../data/weather/weather-test.csv /tmp/weather-posteriors.txt
diff /tmp/weather-posteriors.txt ../data/weather/weather-test.ref.posteriors
\end{lstlisting}

\begin{lstlisting}[style=outputstyle]
ete       chaud  0.609347
hiver     froid  0.727018
mi-saison doux   0.568331
mi-saison doux   0.577112
hiver     froid  0.748287
ete       chaud  0.589928
\end{lstlisting}

\begin{resultbox}
Les probabilités a posteriori calculées sont \textbf{exactement identiques} au fichier
de référence \texttt{weather-test.ref.posteriors}. Les trois fonctions
\texttt{alpha}, \texttt{beta} et \texttt{pgen} sont correctement implémentées.
\end{resultbox}

\subsubsection*{Fichier \texttt{fr-test.word} étiqueté avec posteriors (extrait)}

\begin{lstlisting}[style=outputstyle]
Qu'     PRON   0.996577
est     AUX    0.962172
-ce     PRON   1.000000
qu'     SCONJ  0.843615
une     DET    0.983906
\end{lstlisting}

\begin{notebox}
On observe que les étiquettes à \textbf{forte confiance} (proche de 1) correspondent
aux mots grammaticaux sans ambiguïté (déterminants, ponctuation, auxiliaires).
Les \textbf{confiances basses} ($< 0{,}6$) signalent les cas ambigus, cohérents avec
les erreurs élevées observées pour les étiquettes ADJ, VERB et ADV.
Ces scores de confiance constituent donc un indicateur fiable de la qualité des prédictions.
\end{notebox}

% ═════════════════════════════════════════════════════════════════════════════
\section*{Conclusion}
\addcontentsline{toc}{section}{Conclusion}
% ═════════════════════════════════════════════════════════════════════════════

Ce TP a permis d'implémenter et de valider les composantes essentielles d'un modèle
de Markov caché :

\begin{itemize}
  \item L'\textbf{algorithme de Viterbi} (décodage optimal, 0\% d'erreur sur la météo).
  \item L'\textbf{apprentissage supervisé MLE} (paramètres proches du vrai modèle,
        0\% d'erreur après apprentissage).
  \item L'\textbf{algorithme forward-backward} (posteriors conformes à la référence).
  \item Trois applications : prédiction météo, détection de régions CG en bio-informatique
        (83\% de précision), et étiquetage morpho-syntaxique du français (~92\% de précision).
\end{itemize}

Les principales limites observées concernent le POS-tagging : l'hypothèse de Markov
d'ordre 1 ne capture pas les dépendances longue distance, et le mécanisme \texttt{UNK}
reste une approximation grossière pour les mots inconnus. Des modèles plus avancés
(HMM d'ordre 2, CRF, réseaux de neurones) permettraient d'améliorer ces résultats.

\end{document}
